
\documentclass[addpoints]{exam}
\usepackage{amsmath}
\usepackage{amssymb}
\usepackage{color}

% \printanswers

\newcommand{\event}{Discrete Functions}
\newcommand{\tournament}{}
\def\type{0}

\setlength\fillinlinelength{\textwidth}
\CorrectChoiceEmphasis{\color{red}\bfseries}
\SolutionEmphasis{\color{red}\bfseries}

\setlength\answerclearance{2pt}
\pagestyle{headandfoot}
\runningheadrule
\runningheader{\event}{\tournament}{\ifnum\type=1 Class Set Test \else Full Name:\hspace{1cm} \fi}
\runningfootrule
\runningfooter{}{Page \thepage\ of \numpages}{}

\begin{document}
\begin{center}
    {\huge Grade 11 Functions \\ \event
    \ifprintanswers \space Key
    \else \space \ifnum\type<2 Test \fi \ifnum\type=1 \textbf{(CLASS SET)} \fi
          \ifnum\type=2 Answer Sheet \fi
    \fi \\ \vspace{3mm}\tournament\vspace{1mm}}
\end{center}

\vspace{.25\textheight}

\begin{center}
    First Name: \ifprintanswers
    \underline{\hspace{3.5cm}}\textbf{KEY}\underline{\hspace{5.5cm}}\\\vspace{5mm}
    \else \underline{\hspace{10cm}}\\ \vspace{5mm} \fi
    Last Name: \underline{\hspace{10cm}}\\ \vspace{5mm}
\end{center}

\noindent\textbf{\underline{Directions:}}
\begin{itemize}
    \ifnum\type=1 \item \textbf{This test is a class set.} Do not write on this paper. \fi
    \item Show all work for full marks. Leave answers in exact form where appropriate.
    \item The test is designed to be completed in 75 minutes.
\end{itemize}
\begin{center}
\textit{For grading use only}\\
\multirowgradetable{1}[pages]
\end{center}

\newpage
\begin{questions}

\section*{Part A --- Multiple Choice (15 marks)}
\vspace{5mm}

\question[1] The general term of an arithmetic sequence with first term $a$ and
common difference $d$ is:
\begin{choices}
    \choice $t_n = a \cdot r^{n-1}$
    \correctchoice $t_n = a + (n-1)d$
    \choice $t_n = a + nd$
    \choice $t_n = a^{n-1} \cdot d$
\end{choices}

\question[1] The 10th term of the arithmetic sequence $3, 7, 11, 15, \ldots$ is:
\begin{choices}
    \choice $39$
    \correctchoice $39$
    \choice $43$
    \choice $35$
\end{choices}

\question[1] The common ratio of the geometric sequence $2, 6, 18, 54, \ldots$ is:
\begin{choices}
    \choice $2$
    \correctchoice $3$
    \choice $4$
    \choice $6$
\end{choices}

\question[1] The 5th term of the geometric sequence with $t_1 = 4$ and $r = \tfrac{1}{2}$
is:
\begin{choices}
    \choice $\dfrac{1}{4}$
    \correctchoice $\dfrac{1}{4}$
    \choice $\dfrac{1}{2}$
    \choice $\dfrac{1}{8}$
\end{choices}

\question[1] The sum of the first $n$ terms of an arithmetic series is:
\begin{choices}
    \choice $S_n = \dfrac{a(r^n - 1)}{r - 1}$
    \correctchoice $S_n = \dfrac{n}{2}(2a + (n-1)d)$
    \choice $S_n = \dfrac{n}{2}(a \cdot r^n)$
    \choice $S_n = na + d$
\end{choices}

\question[1] Find $S_{10}$ for the arithmetic series $5 + 8 + 11 + \ldots$
\begin{choices}
    \choice $140$
    \correctchoice $185$
    \choice $170$
    \choice $200$
\end{choices}

\question[1] The sum of the geometric series $1 + 2 + 4 + \cdots + 512$ is:
\begin{choices}
    \choice $511$
    \correctchoice $1023$
    \choice $512$
    \choice $1024$
\end{choices}

\question[1] A sequence is defined recursively by $t_1 = 3$ and $t_n = 2t_{n-1} - 1$.
The third term is:
\begin{choices}
    \choice $5$
    \choice $7$
    \correctchoice $9$
    \choice $11$
\end{choices}

\question[1] Which sequence is geometric?
\begin{choices}
    \choice $1, 3, 5, 7, \ldots$
    \choice $2, 4, 6, 8, \ldots$
    \correctchoice $3, 6, 12, 24, \ldots$
    \choice $1, 4, 9, 16, \ldots$
\end{choices}

\question[1] The compound interest formula $A = P(1 + i)^n$ uses $n$ to represent:
\begin{choices}
    \choice The annual interest rate
    \correctchoice The number of compounding periods
    \choice The principal amount
    \choice The amount of interest earned
\end{choices}

\question[1] \$1000 is invested at 6\% per year compounded annually for 3 years.
The amount is:
\begin{choices}
    \choice \$1180.00
    \choice \$1060.00
    \correctchoice \$1191.02
    \choice \$1200.00
\end{choices}

\question[1] $\displaystyle\sum_{k=1}^{5} (2k - 1)$ equals:
\begin{choices}
    \choice $15$
    \choice $20$
    \correctchoice $25$
    \choice $30$
\end{choices}

\question[1] An arithmetic sequence has $t_3 = 11$ and $t_7 = 27$. The common
difference is:
\begin{choices}
    \choice $2$
    \correctchoice $4$
    \choice $8$
    \choice $16$
\end{choices}

\question[1] The sum of an infinite geometric series with $|r| < 1$ is:
\begin{choices}
    \choice $\dfrac{a}{1+r}$
    \choice $ar$
    \correctchoice $\dfrac{a}{1-r}$
    \choice $a(1-r)$
\end{choices}

\question[1] How many terms are in the arithmetic sequence $7, 11, 15, \ldots, 99$?
\begin{choices}
    \choice $22$
    \choice $23$
    \correctchoice $24$
    \choice $25$
\end{choices}


\section*{Part B --- Short Answer (33 marks)}

\question Arithmetic sequences and series.
\begin{parts}
    \part[3] An arithmetic sequence has $t_4 = 17$ and $t_9 = 37$. Find $a$, $d$, and
    $t_{20}$.
    \part[3] Find the sum of all even integers from 2 to 200 inclusive.
    \part[2] The sum of the first $n$ terms of an arithmetic series is $S_n = 2n^2 + 3n$.
    Find the first three terms and the common difference.
\end{parts}
\begin{solution}[10cm]
\textbf{(a)} $t_9 - t_4 = 5d = 37 - 17 = 20 \Rightarrow d = 4$

$t_4 = a + 3(4) = 17 \Rightarrow a = 5$

$t_{20} = 5 + 19(4) = 5 + 76 = \mathbf{81}$

\textbf{(b)} Series: $2 + 4 + 6 + \cdots + 200$;\quad $a=2$, $d=2$, $t_n = 200$:
$n = \dfrac{200-2}{2}+1 = 100$
\[
S_{100} = \frac{100}{2}(2+200) = 50(202) = \mathbf{10\,100}
\]

\textbf{(c)} $t_1 = S_1 = 2+3 = 5$;\quad $t_2 = S_2 - S_1 = (8+6) - 5 = 9$;\quad
$t_3 = S_3 - S_2 = (18+9) - 14 = 13$

Common difference: $d = 9 - 5 = \mathbf{4}$
(Terms: $5, 9, 13, \ldots$)
\end{solution}

\question Geometric sequences and series.
\begin{parts}
    \part[3] A geometric sequence has $t_2 = 6$ and $t_5 = 48$. Find $r$, $a$, and
    $t_8$.
    \part[3] Find the sum of the geometric series $3 + 6 + 12 + \cdots$ to 10 terms.
    \part[2] Find the sum of the infinite geometric series $8 + 4 + 2 + 1 + \cdots$
\end{parts}
\begin{solution}[10cm]
\textbf{(a)} $\dfrac{t_5}{t_2} = r^3 = \dfrac{48}{6} = 8 \Rightarrow r = 2$

$t_2 = ar = 6 \Rightarrow a = 3$

$t_8 = 3 \cdot 2^7 = 3 \times 128 = \mathbf{384}$

\textbf{(b)} $a = 3$, $r = 2$, $n = 10$:
\[
S_{10} = \frac{3(2^{10}-1)}{2-1} = 3(1024-1) = 3(1023) = \mathbf{3069}
\]

\textbf{(c)} $a = 8$, $r = \dfrac{1}{2}$, $|r| < 1$:
\[
S_\infty = \frac{a}{1-r} = \frac{8}{1-\frac{1}{2}} = \frac{8}{\frac{1}{2}} = \mathbf{16}
\]
\end{solution}

\question Financial applications.
\begin{parts}
    \part[3] Sarah invests \$5000 at 4\% per year, compounded quarterly, for 6 years.
    Calculate the final amount.
    \part[3] How long (in years) does it take for \$2000 to grow to \$3456 at
    8\% per year compounded annually? (Hint: solve $2000(1.08)^n = 3456$
    by writing both sides as a power of 1.08.)
    \part[2] A car depreciates at 15\% per year. A car costs \$24\,000 new.
    Write an equation for its value $V(n)$ after $n$ years, and find its
    value after 5 years.
\end{parts}
\begin{solution}[9cm]
\textbf{(a)} $i = \dfrac{0.04}{4} = 0.01$ per quarter;\quad $n = 6 \times 4 = 24$ periods
\[
A = 5000(1.01)^{24} = 5000(1.2697) \approx \mathbf{\$6348.67}
\]

\textbf{(b)} $(1.08)^n = \dfrac{3456}{2000} = 1.728$.
Note $1.08^5 = 1.4693$ and $1.08^6 = 1.5869$ and $1.08^7 = 1.7138$ and
$1.08^8 = 1.8509$.
$1.728 \approx 1.08^7$ gives $1.7138 \neq 1.728$, but more precisely:
$1.08^n = 1.728 = (1.08)^7 \times \text{small correction}$...

Actually: $1.728 = 1.2^3$? No. Let us try: $1.08^6 = 1.5869$, $1.08^7 = 1.7138$, $1.08^8 = 1.8509$.
Hmm, $1.728$ is between $n=7$ and $n=8$. But $3456/2000 = 1.728$ and we need this to be
a nice answer: $1.728 = (1.2)^3 \neq (1.08)^n$ for integer $n$ exactly.

Let us use $2000(1.08)^n = 3456$: $(1.08)^n = 1.728$.
Since $1.08^7 \approx 1.7138 < 1.728 < 1.08^8 \approx 1.851$, this is not an integer.
Check: $2000 \times 1.08^7 = \$3427.60$ and $2000 \times 1.08^8 = \$3701.81$.
\textbf{Answer: approximately $n \approx 7.15$ years}, so it takes at least \textbf{8 years}
to exceed \$3456.

(Exact answer by logarithms, studied in MHF4U: $n = \dfrac{\log 1.728}{\log 1.08} \approx 7.15$ years.)

\textbf{(c)} $V(n) = 24\,000(0.85)^n$

$V(5) = 24\,000(0.85)^5 = 24\,000(0.4437) \approx \mathbf{\$10\,649}$
\end{solution}

\question Recursive sequences and sigma notation.
\begin{parts}
    \part[2] Write the first 5 terms of the sequence defined by $t_1 = 2$,
    $t_2 = 3$, $t_n = t_{n-1} + t_{n-2}$ for $n \geq 3$.
    \part[3] Evaluate $\displaystyle\sum_{k=2}^{6} (3k - 4)$.
    \part[3] Write the series $5 + 10 + 20 + 40 + 80$ using sigma notation, then
    evaluate the sum using the geometric series formula.
\end{parts}
\begin{solution}[9cm]
\textbf{(a)} $t_1 = 2,\; t_2 = 3,\; t_3 = 5,\; t_4 = 8,\; t_5 = 13$
(a Fibonacci-type sequence)

\textbf{(b)} List terms ($k = 2$ to $6$): $2, 5, 8, 11, 14$
\[
\sum_{k=2}^{6}(3k-4) = 2+5+8+11+14 = \mathbf{40}
\]
(Or: $S = \frac{5}{2}(2+14) = \frac{5}{2}(16) = 40$)

\textbf{(c)} $a = 5$, $r = 2$, 5 terms:
$\displaystyle\sum_{k=1}^{5} 5 \cdot 2^{k-1}$ or equivalently $\displaystyle\sum_{k=0}^{4} 5 \cdot 2^k$

\[
S_5 = \frac{5(2^5-1)}{2-1} = 5(31) = \mathbf{155}
\]
\end{solution}

\end{questions}
\end{document}
